% !TEX root = ../main.tex
\section{Introduction}\label{sec:intro}


With the discovery of a Standard-Model (SM)-like Higgs boson at the LHC in 2012~\cite{ATLAS:2012yve,CMS:2012qbp}, the 
SM is complete. Yet many questions beyond the Standard Model (BSM) remain. One avenue of exploration which has
received considerable attention is new light, feebly interacting particles~\cite{Dedes:2001zia, Feng:2014uja, 
Hewett:2014qja, Alexander:2016aln, Battaglieri:2017aum, Alimena:2019zri, Beacham:2019nyx, Arguelles:2019xgp, 
EuropeanStrategyforParticlePhysicsPreparatoryGroup:2019qin, Agrawal:2021dbo}. Such exotic states are predicted in many
BSM theories and are often long-lived. Theoretical candidates for such long-lived particles (LLPs) range from heavy 
neutral leptons (HNLs), axion-like particles, dark scalars, and dark photons, to electroweakinos in variations of 
supersymmetric models, inelastic dark matter (DM), and many more. See, \textit{e.g.},~Refs.~\cite{Alimena:2019zri, 
Lee:2018pag, Curtin:2018mvb, Knapen:2022afb} for reviews on LLPs. They are usually motivated as explanations of either
the non-vanishing active neutrino masses or of dark matter and could have a spin of $0,\frac{1}{2},1,
\ldots$, and a mass usually ranging from the sub-MeV scale up to the multi-TeV scale. 

As an example, the HNLs (labeled as $N$ in this work) are proposed hypothetical spin-half fermions 
that are SM-singlets which mix with the light active neutrinos. For certain mass values, they can explain 
simultaneously the neutrino masses, the observed dark matter, as well as the baryon asymmetry of the 
Universe~\cite{Asaka:2005pn}. They can give light neutrinos Dirac masses via the Yukawa term $LHN$ in the Lagrangian with $L$ and $H$ being the lepton and Higgs doublets.
This implies unnaturally small Yukawa couplings given the tiny neutrino masses~\cite{deSalas:2020pgw, 
Esteban:2020cvm, Capozzi:2021fjo}.
One can also write down a Majorana mass term in the Lagrangian, leading to light Majorana neutrinos via the seesaw mechanisms~\cite{Minkowski:1977sc, Yanagida:1979as, Gell-Mann:1979vob, Mohapatra:1979ia, Schechter:1980gr}.
While the vanilla type-I seesaw mechanism demands the mixing parameters should be small for the tiny active neutrino masses, larger values of mixing are legitimately conceived in other variations such as the linear seesaw model~\cite{Akhmedov:1995ip, Akhmedov:1995vm, Malinsky:2005bi} and inverse seesaw model~\cite{Mohapatra:1986bd, Gonzalez-Garcia:1988okv}.
Therefore, in phenomenological studies, the HNL mass and the mixing angles with the SM neutrinos are often assumed to be independent parameters. 

Via mixing with the active neutrinos, the HNLs can participate in both charged-current and 
neutral-current interactions, coupled to both gauge bosons directly and the Higgs boson indirectly. 
They can thus be produced from decays of these bosons or from mesons, or through direct production 
at colliders. The HNL can decay leptonically or semi-leptonically, leading to a variety of signatures 
at the different experimental facilities. In particular, GeV-scale HNLs have received substantial 
attention in recent years, for they could originate from rare decays of mesons, which are copiously 
produced, \textit{e.g.}, at beam-dump experiments, $B$-factories, and high-energy hadron colliders. Given the 
strict experimental upper bounds on the mixing of the HNLs and the active neutrinos, the more recent 
focus has been on small mixing angles, for which the GeV-scale HNLs are usually long-lived. Searches 
for these long-lived HNLs have been performed via many different signatures, including searches for 
missing energy, peak searches, as well as searches for displaced vertices (DV). See, 
\textit{e.g.},~Refs.~\cite{hnl_limits, Fernandez-Martinez:2023phj} for summaries of these searches. 
Moreover, one could use the uncertainty on the measurements of the invisible decay width of 
mesons to put upper bounds on the long-lived HNLs, which contribute to the invisible decay width.




Besides the HNLs, supersymmetric electroweakinos, including charginos and neutralinos, are often considered as LLP candidates. See for instance Refs.~\cite{ATLAS:2022rme, CMS:2020atg}.
In particular, a specific type of light neutralino in the GeV mass scale is still allowed by all observational and experimental 
constraints~\cite{Choudhury:1999tn, Gogoladze:2002xp, Dreiner:2009er, Dreiner:2009ic, Grifols:1988fw, Ellis:1988aa, 
Lau:1993vf, Dreiner:2003wh, Dreiner:2006gu, Profumo:2008yg, Dreiner:2011fp, Dreiner:2013tja}; they are necessarily 
bino-like~\cite{Dreiner:2009ic,Choudhury:1999tn} and have to decay to avoid overclosing the Universe~\cite{Hooper:2002nq,
Calibbi:2013poa, Bechtle:2015nua}. One possibility is to consider R-parity-violating supersymmetry (RPV-SUSY) (see 
Refs.~\cite{Dreiner:1997uz,Barbier:2004ez,Mohapatra:2015fua} for reviews), where the light binos decay via small but 
non-vanishing RPV couplings (see Ref.~\cite{Domingo:2022emr} for a detailed study of light bino decays).
The minimal version is known as the R-parity-violating Minimal Supersymmetric Standard Model (RPV-MSSM) \cite{Allanach:2003eb}. The 
RPV-MSSM solves the SM hierarchy problem as in the MSSM, and also predicts a very rich phenomenology at 
colliders~\cite{Dreiner:1991pe, deCampos:2007bn, Dercks:2017lfq, Dreiner:2023bvs}.
\textit{A priori} it is unknown if R-parity is conserved or broken, SUSY models with R-parity conservation or violation are equally legitimate, \textit{e.g.}, see 
Refs.~\cite{Ibanez:1991pr, Dreiner:2005rd}.
Moreover, the RPV-MSSM can explain several experimental anomalies reported in recent years including the $B$-meson anomalies~\cite{Trifinopoulos:2019lyo, Hu:2020yvs, Hu:2019ahp, Zheng:2021wnu, 
BhupalDev:2021ipu}, the ANITA anomaly~\cite{Altmannshofer:2020axr, Collins:2018jpg}, as well as the anomalous magnetic
moment of the muon~\cite{Hu:2019ahp, Zheng:2021wnu, BhupalDev:2021ipu}. If R-parity is broken, one can write down 
operators which violate baryon- or lepton-number.
Allowing all these operators to be non-vanishing would lead to a too fast proton decay rate, in conflict with the current experimental measurements~\cite{Chamoun:2020aft, ParticleDataGroup:2022pth}, unless all their couplings are extremely small.
Therefore, we assume the discrete anomaly-free baryon triality symmetry $B_3$ \cite{Dreiner:2006xw, Lee:2010vj, Dreiner:2011ft}, so that baryon-number is conserved.
In this work, we restrict ourselves to the lepton-number-violating terms only.
Further, the light bino is the lightest supersymmetric particle (LSP) in our study and decays only into SM particles via RPV couplings.

Via the lepton-number-violating RPV operators, the light bino decays lead to very similar final states as the NHL 
decays.
Moreover, the corresponding RPV couplings are all bounded to be small by various low-energy observables and
collider searches~\cite{Barbier:2004ez, Allanach:1999ic, Bolton:2021hje, Bansal:2018dge, Bansal:2019zak}. The GeV-scale
binos are hence expected to be long-lived, too, resulting in signatures such as missing energy and displaced vertices at various experiments.
These similarities raise the question: is it possible to recast the extensive exclusion bounds 
on the HNLs in the literature into corresponding bounds on the light binos in the RPV-SUSY?
In this work, we answer this question positively, by compiling a list of bounds on long-lived HNLs obtained in searches for all types of signatures 
mentioned above and recasting them into exclusion limits on the RPV-SUSY couplings as functions of the light bino mass for 
a selected list of benchmark scenarios.\footnote{Recently during the completion of this work, Ref.~\cite{Fernandez-Martinez:2023phj} appeared on arXiv; it employed similar strategies to recast the bounds on the HNLs in the minimal scenarios into those on the HNLs in effective field 
theories.} We focus on exclusion bounds acquired in past experiments, as well as predicted search sensitivities for experiments that are ongoing 
or under construction. For future (and not yet approved) experiments, we consider only 
\texttt{MoEDAL-MAPP2}~\cite{Pinfold:2019nqj, Pinfold:2019zwp} and \texttt{FASER2}~\cite{FASER:2018eoc} with 300 and 
3000\,fb$^{-1}$ integrated luminosity, respectively, as they would be the successors of some ongoing experiments at the 
LHC, while the other future concepts such as \texttt{MATHUSLA}~\cite{Curtin:2018mvb, Chou:2016lxi, MATHUSLA:2020uve}, and \texttt{ANUBIS}~\cite{Bauer:2019vqk} are independent ones and are hence not studied here. We do not consider the approved experiment \texttt{Hyper-Kamiokande}~\cite{Abe:2011ts, Hyper-KamiokandeWorkingGroup:2013hcb} for there is no 
available HNL-search sensitivity prediction that can be used with our recasting methods. 





In the following section, we give the model basics of light binos in the RPV-SUSY and of the HNLs that mix with active neutrinos.
The considered experiments are introduced in Sec~\ref{sec:exp}, along with an explanation of our recasting procedure.
We then present our numerical results for some representative benchmark scenarios in Sec.~\ref{sec:results}.
Finally, in Sec.~\ref{sec:conclusions} we conclude the paper with a summary.




